% Shared formatting for course notes. Loaded by each course's main.tex.
\usepackage[T1]{fontenc}
\usepackage[utf8]{inputenc}
\usepackage{lmodern}
\linespread{1.15}
\usepackage[margin=1in]{geometry}
\usepackage{amsmath,amssymb,amsthm}
\usepackage{mathtools}
\usepackage{graphicx}
\usepackage{booktabs}
\usepackage{enumitem}
\usepackage{xcolor}
\usepackage{hyperref}
\usepackage{tikz}
\usetikzlibrary{decorations.pathreplacing,positioning}
\usepackage{xparse}
\usepackage{etoolbox}

\hypersetup{colorlinks=true,linkcolor=blue!50!black,urlcolor=blue!50!black}
\setlength{\parindent}{0pt}
\setlength{\parskip}{0.5em}
\setlist{nosep}

% Unnumbered section headings: entries show as "Week 1", not "1 Week 1",
% both in the body and in the table of contents.
\setcounter{secnumdepth}{0}

% Each \section (i.e. each week's entry) starts on its own new page. The
% \newpage that used to follow \tableofcontents in each course's main.tex
% is redundant now and has been removed, so week 1 isn't pushed to a blank
% page by this.
\pretocmd{\section}{\newpage}{}{\errmessage{Failed to patch \string\section}}

% Tracks whether \section was the most recent sectioning/lecture command,
% so a \newlecture immediately following a \section (the usual "week's
% first lecture" case) doesn't also throw in its own \newpage, which would
% otherwise leave a blank page between the two. Set true by \section, and
% cleared by \newlecture whether or not it was true, so it only fires for
% a \newlecture that directly follows a \section (nothing else in between).
\newbool{newlecture@justsection}
\pretocmd{\section}{\booltrue{newlecture@justsection}}{}{\errmessage{Failed to patch \string\section}}

% Custom head format: bold, unparenthesized title, with a \quad gap after
% the number (e.g. "Theorem 3.2  Distribution" instead of "Theorem 3.2
% (Distribution)."), reusing amsthm's usual spacing/punctuation otherwise.
% When a title is given, the body text starts on its own new line below the
% head instead of running inline after it; untitled (auto-numbered-only)
% instances stay inline as usual. Body text is upright (not italic) for all
% six environments. Two variants share this head format but differ in
% end-of-head punctuation: theorem/lemma/proposition/exercise keep a
% trailing period, definition/example drop it (no period after the title).
\makeatletter
\newcommand{\thmpunct}{.}
\newtheoremstyle{boldhead}{}{}{\normalfont}{}{\bfseries}{\thmpunct}{5pt plus 1pt minus 1pt}%
    {\thmname{#1}\thmnumber{\@ifnotempty{#1}{ }#2}\thmnote{\@ifnotempty{#3}{\quad#3}}%
     \@ifnotempty{#3}{\def\thmheadnl{\newline}\thm@headsep\z@skip}{}}
\newtheoremstyle{boldheadnp}{}{}{\normalfont}{}{\bfseries}{}{5pt plus 1pt minus 1pt}%
    {\thmname{#1}\thmnumber{\@ifnotempty{#1}{ }#2}\thmnote{\@ifnotempty{#3}{\quad#3}}%
     \@ifnotempty{#3}{\def\thmheadnl{\newline}\thm@headsep\z@skip}{}}
\makeatother

\theoremstyle{boldhead}
\newtheorem{theorem}{Theorem}
\newtheorem{lemma}[theorem]{Lemma}
\newtheorem{proposition}[theorem]{Proposition}
\newtheorem{corollary}[theorem]{Corollary}
\newtheorem{exercise}[theorem]{Exercise}
\theoremstyle{boldheadnp}
\newtheorem{definition}[theorem]{Def.}
\newtheorem{example}[theorem]{Example}
\theoremstyle{remark}
\newtheorem*{remark}{Remark}

% \begin{theorem}[Title][2.7]: both optional arguments are new on top of
% the environments' usual single bracket. The first argument still behaves
% exactly like plain amsthm's optional title (existing notes using
% \begin{example}[Some Name] etc. are unaffected). The second argument, if
% given, overrides the displayed number (e.g. to match a textbook's own
% numbering) instead of the shared theorem counter's automatic value; the
% counter itself still advances normally, so later auto-numbered
% environments are unaffected too.
\newcommand{\addtitleandnumber}[1]{%
    \csletcs{#1orig}{#1}%
    \csletcs{end#1orig}{end#1}%
    \RenewDocumentEnvironment{#1}{o o}{%
        \IfValueT{##2}{\ifblank{##2}{}{\renewcommand{\thetheorem}{##2}}}%
        \IfValueTF{##1}{\ifblank{##1}{\csname#1orig\endcsname}{\csname#1orig\endcsname[##1]}}{\csname#1orig\endcsname}%
    }{\csname end#1orig\endcsname}%
}
\addtitleandnumber{theorem}
\addtitleandnumber{lemma}
\addtitleandnumber{proposition}
\addtitleandnumber{corollary}
\addtitleandnumber{definition}
\addtitleandnumber{example}
\addtitleandnumber{exercise}

% Common mathematical notation.
\newcommand{\R}{\mathbb{R}}
\newcommand{\N}{\mathbb{N}}
\newcommand{\Z}{\mathbb{Z}}
\newcommand{\E}{\mathbb{E}}
\newcommand{\Prob}{\mathbb{P}}
\DeclareMathOperator{\Var}{Var}
\DeclareMathOperator{\Cov}{Cov}

% \newlecture[title][date]: bold, centered marker for the start of a new
% lecture. Both arguments are optional. Numbering runs continuously across
% the whole document (i.e. across all weeks' entries, not reset per
% \section), so the Nth \newlecture in the course is always "Lecture N".
% Each lecture starts on its own new page, except when it directly follows
% a \section heading (which already started a new page of its own), so the
% two stay together instead of leaving a blank page in between.
\newcounter{lecture}
\renewcommand{\thelecture}{\arabic{lecture}}
\NewDocumentCommand{\newlecture}{o o}{%
    \stepcounter{lecture}%
    \ifbool{newlecture@justsection}{}{\newpage}%
    \boolfalse{newlecture@justsection}%
    \par\vspace{1em}
    \begin{center}
        \textbf{\Large Lecture \thelecture\IfValueT{#1}{: #1}}\\[2pt]
        \IfValueT{#2}{\normalsize #2}
    \end{center}
    \vspace{0.5em}%
}

% Handwritten-note-style helpers: red definition labels, blue margin
% annotations, yellow highlights, underlined mini-headers.
\newcommand{\deflabel}[2]{\par\noindent{\color{red!70!black}\textbf{Def #1}}\quad\textbf{#2}\par\nobreak}
\newcommand{\exlabel}[2]{\par\noindent\textbf{Example #1}\quad\textbf{#2}\par\nobreak}
\newcommand{\heading}[1]{\par\noindent{\color{red!70!black}\textbf{#1}}\par}
\newcommand{\header}[1]{\par\vspace{0.5em}\noindent\textbf{#1}\par\vspace{-0.2em}}
\newcommand{\aside}[1]{\ifmmode{\color{blue!55!black}\text{\itshape #1}}\else{\color{blue!55!black}\itshape #1}\fi}
\newcommand{\hly}[1]{\colorbox{yellow!55}{$#1$}}
\newcommand{\hlt}[1]{\colorbox{yellow!55}{#1}}
% \boldmath makes any math the argument switches into (e.g. \rb{$r$}) use
% the bold math version automatically, so callers don't need to manually
% wrap math content in \boldsymbol{}.
\newcommand{\rb}[1]{{\color{red!70!black}\boldmath\textbf{#1}}}
\renewcommand{\b}[1]{\textbf{#1}}
\definecolor{navy}{RGB}{0,0,128}
\newcommand{\dbb}[1]{{\color{navy}\textbf{#1}}}
\newcommand{\db}[1]{{\color{navy}#1}}
